\section{Fixed input laws and independent shears}
\label{shear:sec:directions}

The local proof obtained balance by sampling one omitted block.
For an arbitrary fixed input list, a previous switch supplies
balance instead: it crosses the two halves relevant to the next
direction, and the intervening switches preserve those halves.
This gives an ordinary tuple law for every fixed list. We then
record an independent shear realization of fresh directions
and several estimates that retain an average over inputs.
The realization will also control which auxiliary data may be
disclosed in the later history arguments.

We use $N=n=2^d$ and the coordinate process
$X^{\boldsymbol\iota}$ of Lemma~\ref{lem:frames}, where
$\boldsymbol\iota=(\iota_1,\ldots,\iota_d)$ is any ordering of
the standard axes and $i(s)=\iota_{1+((s-1)\bmod d)}$. Thus
\begin{equation}\label{shear:eq:coordinate-chain}
 X^{\boldsymbol\iota}_0=\id,\qquad
 X^{\boldsymbol\iota}_s=h_sX^{\boldsymbol\iota}_{s-1},\qquad
 h_s\sim U_{C_{e_{i(s)}}}
\end{equation}
with independent increments. The order $(1,\ldots,d)$ gives the
rotating-frame chain $X$; its physical permutation is $R^sX_s$.
Other starting phases and the reversed order are included by
choosing another $\boldsymbol\iota$.

\subsection{Balance for every fixed list}

We will use the following elementary concentration estimates.
\begin{lemma}[Concentration estimates]\label{lem:concentration}
Let $X_1,\ldots,X_r$ be independent random variables in $[0,1]$
and put $S=\sum_iX_i$. For $a\ge0$ and $r\ge1$,
\[
 \Pr(S-\E S\ge a),\quad\Pr(S-\E S\le-a)\le e^{-2a^2/r}.
\]
More generally, if $(M_j)_{j=0}^r$ is a real martingale whose
increments have absolute value at most $c_j$, either tail at
distance $a$ is at most
$\exp[-a^2/(2\sum_jc_j^2)]$ when the denominator is positive.
If all $c_j$ vanish, the martingale is constant. In particular,
if $X$ counts the points of a uniform $m$-subset in a fixed half
of a finite set of even size, then
\[
 \Pr(X<m/4)\le e^{-m/32}\qquad(m\ge1).
\]
\end{lemma}
\begin{proof}
For a variable $Z$ in an interval of length $b$, put
$\psi(s)=\log\E e^{sZ}$. Its second derivative is the variance
under the exponentially tilted law, at most $b^2/4$. Integrating
from $\psi(0)=0$ and $\psi'(0)=\E Z$ gives
$\E e^{s(Z-\E Z)}\le e^{s^2b^2/8}$. Multiplication and
minimization in $s$ prove the independent-variable bound. The
conditional version, with interval length $2c_j$, gives
$\E e^{s(M_r-M_0)}\le e^{s^2\sum_jc_j^2/2}$ and the martingale
bound. These are the exponential-moment arguments underlying the
inequalities of Hoeffding and Azuma
\cite{hoeffding1963,azuma1967}.

Expose a uniform subset in a uniform random order and form the
Doob martingale for its final half-count. The expected final count
changes by at most one at each draw: completions following two
possible next draws can be coupled by exchanging those points.
The initial expectation is $m/2$. Use $r=m$, $c_j=1$, and
$a=m/4$.
\end{proof}

For the fixed-list estimate, generate $Z$ directly as in
Section~\ref{sec:marginals}: its first $d$ directions are
$v_j=e_{\iota_j}$, and for $s>d$ choose $v_s$ uniformly outside
the hyperplane
\begin{equation}\label{shear:eq:hyperplane}
 W_s=\Span(v_{s-d+1},\ldots,v_{s-1}).
\end{equation}
Given the complete schedule, the switches are independent and
uniform. Every $d$ consecutive directions form a basis. For a
fixed input list, take deterministic initial data $\Xi$ in
\eqref{eq:marginal-filtration}.

\begin{lemma}[Balance after the previous transverse switch]
\label{shear:lem:fixed-balance}
Fix any $\ell<N$ observed labels, and put $m=N-\ell$. For $s>d$,
let $D_{s-1}$ be their available positions just before step $s$.
Then
\begin{equation}\label{shear:eq:fixed-balance}
 \Pr\{\min(|D_{s-1}\cap W_s|,|D_{s-1}\setminus W_s|)<m/4\}
 \le2e^{-m/8}.
\end{equation}
The probability averages the observed paths; it is not
conditional on them.
\end{lemma}
\begin{proof}
Fix the complete direction schedule. The direction $v_{s-d}$
crosses the two cosets of $W_s$, because the $d$ directions
starting at $s-d$ form a basis. The intervening directions lie
in $W_s$ and preserve each coset.

Condition on the observed configuration immediately before step
$s-d$. Let $a$ be the number of matching pairs with two available
positions and $b$ the number with one, so $m=2a+b$. Immediately
after this crossing switch, the available count on either chosen
side is $a+\operatorname{Bin}(b,1/2)$, using $b$ independent fair
coins. Its mean is $m/2$. If $b>0$, Lemma~\ref{lem:concentration}
bounds a deviation of at least $m/4$ by
$2\exp(-2(m/4)^2/b)\le2e^{-m/8}$; if $b=0$ the count is exactly
$m/2$. The intervening switches preserve the count. Remove the
conditioning to obtain the claim.
\end{proof}

\begin{proposition}[Uniform large-marginal estimates]
\label{shear:prop:fixed-marginal}
Let $L\ge2$ be a fixed integer dividing $N$, and let
$0\le q\le N-N/L$. For every ordered list $C$ of $q$ distinct
labels, every periodic ordering $\boldsymbol\iota$ of the standard
axes, and every integer $T\ge d$,
\begin{equation}\label{shear:eq:general-fixed-marginal}
 \|\mathcal L(X^{\boldsymbol\iota}_T(C))-\pi_q\|_{\TV}
 \le\frac{N^{3/2}}2
 \left(e^{-(T-d)/(4L)}+2e^{-N/(8L)}\right)^{1/2}.
\end{equation}
The same bound holds for $(X^{\boldsymbol\iota}_T)^{-1}(C)$.
For the physical shuffle and its inverse, the corresponding
position laws satisfy the same bound.

The following choices used below give distance at most
$N^{3/2}\sqrt{a_T}/2$:
\begin{align}
 L=8,\quad T=256d:&\quad
 a_T=(31/32)^{255d}+64e^{-N/64},
 \label{shear:eq:256-marginal}\\
 L=8,\quad T=257d:&\quad
 a_T=e^{-8d}+2e^{-N/64},
 \label{shear:eq:257-marginal}\\
 L=32,\quad T=2049d:&\quad
 a_T=e^{-16d}+256e^{-N/256}.
 \label{shear:eq:2049-marginal}
\end{align}
Each of these three substituted bounds tends to zero as
$d\to\infty$, uniformly in the list and axis order. The last
energy bound is $o(N^{-4})$.
\end{proposition}
\begin{proof}
For a successive hidden-card problem, its number $m$ of
available positions is at least $N/L$. Off the event in
\eqref{shear:eq:fixed-balance}, Lemma~\ref{shear:lem:energy}
gives conditional expected loss at least
$\gamma\Delta_t$, where $\gamma=1/(4L)$. The loss is always
nonnegative. Equation~\eqref{eq:marginal-balance-recurrence}
therefore gives, for $t\ge d$,
\begin{equation}\label{shear:eq:fixed-recursion}
 \E\Delta_{t+1}\le(1-\gamma)\E\Delta_t
       +2\gamma e^{-N/(8L)}.
\end{equation}
Starting from $\Delta_d\le1$ yields
$\E\Delta_T\le(1-\gamma)^{T-d}+2e^{-N/(8L)}$.
Apply Lemma~\ref{shear:lem:tuple} with terminal data
$\Lambda=\mathbf v$. For every complete schedule, the first
basis is $(e_{\iota_1},\ldots,e_{\iota_d})$, so $L_0=I$ in
Lemma~\ref{lem:frames}. The same-list equality
\eqref{eq:frames:tuple-input-map} transfers the conditional
tuple bound to $X^{\boldsymbol\iota}$, giving
\eqref{shear:eq:general-fixed-marginal}.

Each switch is an involution. Reversing the independent factors
therefore identifies $(X^{\boldsymbol\iota}_T)^{-1}$ with a
coordinate chain whose first axis is $i(T)$ and whose periodic
order is reversed. This is another permitted ordering, so the
same estimate applies. The relation $Y_T=R^TX_T$ transfers the
forward physical law by an output bijection. Its inverse is
$X_T^{-1}R^{-T}$; the input bijection only changes the fixed
list, and the bound is uniform in that list.

For \eqref{shear:eq:256-marginal}, weaken the additive term in
\eqref{shear:eq:fixed-recursion} to $2e^{-N/64}$ and sum the
geometric series, giving the factor $64$. At $257d$ keep the
sharper recurrence and use $(31/32)^{256d}\le e^{-8d}$.
More generally, at $T=(1+64L)d$, the weakened additive term
$2e^{-N/(8L)}$ gives $e^{-16d}+8Le^{-N/(8L)}$, yielding
\eqref{shear:eq:2049-marginal}. These rates dominate the factor
$N^3$ in the squared tuple bound, and $e^{-16d}=o(N^{-4})$.
\end{proof}

\subsection{An independent shear realization}

The fixed-list proof above uses the directly generated direction
process. The next construction couples a process with a triangular
first basis and the same later uniform-complement rule to a
coordinate chain with the same fixed inputs. It also proves the
conditional path identity needed when the auxiliary shear entries
are retained at terminal time.

\begin{lemma}[Independent shear realization]\label{shear:lem:keys}
Fix an axis order $\boldsymbol\iota$ and a horizon $T\ge d$. Independently
of the increments in \eqref{shear:eq:coordinate-chain}, choose
independent fair bits $b_{sj}$ for $j\ne i(s)$, put
$b_{s,i(s)}=0$, and write $\Lambda$ for all these entries. Set
\[
 S_s=I+e_{i(s)}b_s,\qquad
 A_0=I,\quad A_s=A_{s-1}S_s,\qquad
 v_s=A_{s-1}e_{i(s)},\qquad
 Z_s=A_sX^{\boldsymbol\iota}_s.
\]
Conditional on $\Lambda$, the increments of $Z$ are independent
uniform switches in directions $v_s$. More precisely, if
$\mathsf P_{v_1,\ldots,v_t}$ is the product-switch path law
through time $t$ starting at the identity, then
\begin{equation}\label{shear:eq:prefix-factorization}
 \mathcal L((Z_s)_{s=0}^t\mid\Lambda)
 =\mathsf P_{v_1,\ldots,v_t}\qquad(0\le t\le T).
\end{equation}
Every $d$ consecutive directions form a basis. For $s>d$,
conditional on $v_1,\ldots,v_{s-1}$, the next direction is
uniform on $V\setminus W_s$, with $W_s$ as in
\eqref{shear:eq:hyperplane}. This law is unchanged by additionally
observing any part of the virtual path prefix through time $s-1$.
\end{lemma}
\begin{proof}
The map $S_s$ is an involution in $C_{e_{i(s)}}$, because
$b_se_{i(s)}=0$. The increment of $Z$ is
\[
 A_{s-1}(S_sh_s)A_{s-1}^{-1}.
\]
For fixed $\Lambda$, the factors $S_sh_s$ are independent and
uniform in their coordinate matching groups. Conjugation gives
the asserted product-switch law, including the joint prefix
identity \eqref{shear:eq:prefix-factorization}.

Extend $i(s)$ periodically beyond $T$ and define the deterministic
next-read time of a column by
\[
 \tau(s,j)=\min\{t>s:i(t)=j\}\qquad(j\ne i(s)).
\]
Right multiplication by $S_s$ leaves the selected column
unchanged and adds $b_{sj}v_s$ to column $j$. Thus the direction
at a column's first read, for $t\le d$, is
\[
 v_t=e_{i(t)}+\sum_{1\le u<t}b_{u,i(t)}v_u.
\]
These first $d$ directions form a triangular basis. At later
reads,
\begin{equation}\label{shear:eq:key-recurrence}
 v_t=v_{t-d}+\sum_{t-d<u<t}b_{u,i(t)}v_u\qquad(t>d).
\end{equation}
In either display the coefficients used at time $t$ are exactly
the entries whose next-read time is $t$. Inductively, all earlier
directions are functions only of entries with next-read time
less than $t$. The batch read at $t$ is therefore independent
and fair conditional on those directions. The recurrence also
shows inductively that every consecutive $d$ directions form a
basis. For $t>d$, its fresh $d-1$ coefficients make $v_t$
uniform on
$v_{t-d}+W_t=V\setminus W_t$.

Finally, \eqref{shear:eq:prefix-factorization} says that the
entire virtual path prefix and $\Lambda$ are conditionally
independent given the direction prefix. Hence observing any
portion of that path prefix cannot reveal the fresh batch that
determines the next direction.
\end{proof}

For initial data $\Xi$ independent of $(\Lambda,(Z_s)_{s=0}^T)$
and determining the observed list, the joint prefix factorization
gives, for every $0\le t\le T$, $1\le j\le q$, and $x\in V$,
\begin{equation}\label{shear:eq:terminal-factorization}
 \Pr\{Z_t(c_j)=x\mid\Lambda,\Xi,\,
              Z_s(c_i):0\le s\le t,\ i<j\}
 =p_t^j(x),
\end{equation}
where $p_t^j$ uses the chronological field
\eqref{eq:marginal-filtration}. At $t=T$ this is precisely the
terminal identification required by Lemma~\ref{shear:lem:tuple}.
For $t<T$ it concerns only predecessor path prefixes through
$t$; conditioning on their future paths requires the separate
path-cylinder argument in Section~\ref{frames:chapter}.
The next direction is averaged in the smaller chronological field
before $\Lambda$ is disclosed.

The coupling also preserves every fixed input list conditionally
on the shears:
\begin{equation}\label{shear:eq:fixed-list-frame}
 \|\mathcal L(Z_T(C)\mid\Lambda)-\pi_q\|_{\TV}
 =\|\mathcal L(X^{\boldsymbol\iota}_T(C))-\pi_q\|_{\TV}.
\end{equation}
Indeed, $X^{\boldsymbol\iota}$ is independent of $\Lambda$ and
$Z_T=A_TX^{\boldsymbol\iota}_T$, with $A_T$ a bijection on
outputs. Here $A_0=I$ although the first directions are generally
triangular combinations of the coordinate axes. This same-input
identity is a feature of the shear realization; the arbitrary
first-basis statement in Lemma~\ref{lem:frames} instead has the
input map $L_0^{-1}$.

\begin{remark}[Equivalent inverse-frame factorization]
\label{shear:rem:factorization}
The same coupling has a useful inverse notation. Let $Q_s$ be
the original independent coordinate increments and choose the
independent shears $S_s$ above. Put
$B_0=I$ and $B_s=S_sB_{s-1}=A_s^{-1}$. Define
\begin{equation}\label{shear:eq:inverse-coupling}
 M_s=B_{s-1}^{-1}(S_sQ_s)B_{s-1}
     =B_s^{-1}Q_sB_{s-1},\qquad
 X^{\boldsymbol\iota}_T=B_TM_T\cdots M_1.
\end{equation}
The second identity follows by telescoping. Conditional on all
$S_s$, the factors $S_sQ_s$ are independent uniform coordinate
switches, so $M_s$ has virtual direction
$B_{s-1}^{-1}e_{i(s)}$. This is the same virtual process and the
same prefix and terminal factorization as above. The map $B_T$
acts only on outputs, leaving a fixed input list unchanged.
The posterior identity concerns virtual positions; once $B_T$
is disclosed, both the virtual endpoint law and its uniform
comparison may be pushed through that known output map.
\end{remark}

The fixed-list balance has a second interpretation in this
coupling. For $s>d$, the current columns of $A_{s-1}$ other
than column $i(s)$ are triangular combinations of the last
$d-1$ selected columns, with diagonal coefficients one. Hence
\[
 W_s=A_{s-1}\{x:x_{i(s)}=0\}.
\]
The two available counts in \eqref{shear:eq:fixed-balance} are
therefore the counts according to bit $i(s)$ in the coordinate
chain. That bit was last refreshed at step $s-d$ and is
unchanged by the intervening coordinate layers. The same
$a+\operatorname{Bin}(b,1/2)$ calculation proves the balance
directly in the original coordinate frame.

\section{Estimates that average the input labels}
\label{shear:sec:averaged}

Sampling the inputs supplies balance without using the previous
transverse switch. These estimates average the distance for each
specified input list; they do not replace that family of
conditional distances by the distance of a mixture. This
distinction permits a later choice of one deterministic partition
or the construction of kernels for all starting states.

\subsection{A reservoir of omitted labels}

\begin{proposition}[A reservoir of omitted labels]
\label{shear:prop:reservoir}
Let $d\ge3$, $b=N/8$, and $T=257d$. Let $B$ be a uniform
$b$-subset of input labels, independent of the shuffle. For any
fixed ordering of $V\setminus B$, let $\delta_T(B)$ be the
distance of its image tuple under $X_T$ from uniform distinct
positions. With $\theta=2^{1/4}(3/4)<1$,
\begin{equation}\label{shear:eq:reservoir-marginal}
 \E_B\delta_T(B)
 \le\frac{N^{3/2}}2
 \left(e^{-8d}+2(N+1)\theta^{N/8}\right)^{1/2}=o(1).
\end{equation}
\end{proposition}
\begin{proof}
Use the direct direction process with standard first basis,
sampled independently of $B$, and take $\Xi=B$ in the online construction. At every time the
available set for a hidden card contains $Z_t(B)$. Fix the full
direction schedule and the full permutation $Z_t$, but not $B$.
The set $Z_t(B)$ is then a uniform $b$-subset, and each coset of
the next-direction hyperplane has size $N/2$.

Its count in a specified half has the distribution of the number
of ones among the first $b$ of $N$ independent fair bits,
conditional on their total being $N/2$. This conditioning event
has probability at least $1/(N+1)$, because the central binomial
coefficient is maximal. For the unconditioned first-$b$ count $J$,
\[
 \Pr\{J\le b/4\}
 \le2^{b/4}\E2^{-J}
 =\bigl(2^{1/4}3/4\bigr)^b.
\]
The inclusion of $Z_t(B)$ in the available set and a union
bound give
$p_{\mathrm{bad}}=2(N+1)\theta^b$ as an upper bound for the
probability that either available half has fewer than $b/4$
positions. This bound is unconditional in the observed paths.

Off this event the conditional expected loss is at least
$\Delta_t/32$. Equations \eqref{eq:marginal-balance-recurrence}
and \eqref{eq:marginal-cross-energy} give
\[
 \E\Delta_{t+1}\le(31/32)\E\Delta_t+p_{\mathrm{bad}}/32,
 \qquad
 \E\Delta_T\le e^{-8d}+p_{\mathrm{bad}}.
\]
Use Lemma~\ref{shear:lem:tuple} with $\Lambda=\mathbf v$.
For each fixed schedule and $B$, the standard first basis gives
$L_0=I$, so the conditional auxiliary distance equals the
distance for $X_T$ on that same complement. Averaging over $B$
proves the bound.
\end{proof}

\subsection{Uniform input tuples after a deterministic first sweep}

\begin{lemma}[Averaged tuple bound]
\label{lifts:averaged-tuples}
Fix an integer $b\ge2$ dividing $N$, put $m=N/b$, and let
$t\ge d$ be an integer. Let $C$ be a uniformly sampled ordered
tuple of $N-m$ distinct input labels, independent of the
shuffle, and let $\mu_C$ be its time-$t$ physical image law with
$C$ fixed. Then
\begin{equation}\label{lifts:averaged-bound}
 \E_C\|\mu_C-\pi_{N-m}\|_{\TV}
 \le\frac{N^{3/2}}2
 \left((1-1/(4b))^{t-d}+2\vartheta^{N/b}\right)^{1/2},
 \qquad \vartheta=\frac32\,2^{-3/4}<1.
\end{equation}
The same bound holds for a fixed ordered list of $N-m$ labels
when its conditional distance is averaged over a uniformly
chosen starting deck.
\end{lemma}
\begin{proof}
Use the direct process with the standard first basis and take
$\Xi=C$. For a hidden label at index $j$, let $h=N-j+1\ge m$ be
the number of available positions. Fix the complete direction
and switch history, but not $C$. At each fixed time, the
preceding positions form a uniform ordered sample, so the
available set is a uniform $h$-subset. The next-direction
hyperplane is fixed under this conditioning.

If $X$ is the count of a uniform $h$-subset in a fixed half,
expand the product of $1+I_i$ over an ordered sample without
replacement. Any specified $\ell$ sample indices lie in that
half with probability at most $2^{-\ell}$. Thus
\begin{equation}\label{shear:eq:uniform-subset-moment}
 \E2^X\le(3/2)^h,\qquad
 \Pr\{\min(X,h-X)<h/4\}\le2\vartheta^h.
\end{equation}
The second inequality follows by applying Markov's inequality
to the count in each half above $3h/4$. It is used only after
averaging out the input list, not conditional on observed paths.

On the balanced event, the conditional expected loss is at
least $h\Delta_s/(4N)$. With
$\gamma_j=h/(4N)$, the unconditional recurrence and its
iteration give
\[
 \E\Delta_{s+1}\le(1-\gamma_j)\E\Delta_s
                   +2\gamma_j\vartheta^h,\qquad
 \E\Delta_t\le(1-1/(4b))^{t-d}+2\vartheta^m.
\]
Lemma~\ref{shear:lem:tuple}, with $\Lambda=\mathbf v$, bounds
the average distance conditional on the full schedule and list.
For every fixed schedule the first basis is standard, so the
same-list frame equality transfers that distance to $X_t$;
the terminal rotation transfers it to the physical shuffle.
This proves \eqref{lifts:averaged-bound}. A uniform starting
deck sends any fixed ordered label list to a uniform input
tuple, which proves the last assertion.
\end{proof}

\begin{corollary}[Block lengths for the averaged input]
\label{lifts:tuple-constants}
For $b=16$, the average distance in
Lemma~\ref{lifts:averaged-tuples} at $t=513d$ is at most
\[
 \varepsilon_N=\frac12N^{3/2}
       \left(N^{-8}+2\vartheta^{N/16}\right)^{1/2}=o(1).
\]
At $t=2049d$ it is at most $N^{-10}$ for sufficiently large
$d$. For $b=1024$ and $d\ge10$, time $t=(1+32b)d$ gives average distance
$o(1)$.
\end{corollary}
\begin{proof}
The respective contracting times are $512d$, $2048d$, and
$32bd$, with energy factors at most $e^{-8d}$, $e^{-32d}$,
and $e^{-8d}$. Use $e^{-8d}\le N^{-8}$ for the first and
$N^{3/2}e^{-16d}=o(N^{-10})$ for the second. Exponential
decay in $N$ handles the other term in each case.
\end{proof}

There is also a concentration proof of the last assertion that
uses the subset martingale in Lemma~\ref{lem:concentration}.
A union bound over the two halves gives imbalance probability
at most $2e^{-h/32}$. With $c_0=1/32$, the weaker recurrence
\begin{equation}\label{lifts:azuma-recurrence}
 \E\Delta_{s+1}\le(1-1/(4b))\E\Delta_s+2e^{-c_0N/b}
\end{equation}
gives $\E\Delta_t\le e^{-8d}+8b e^{-c_0N/b}$ at
$t=(1+32b)d$. The expected endpoint error of one conditional
card law is $o(1/N)$ and its sum is $o(1)$. This is an
alternative proof of the averaged input used later for omitted
sets.

\subsection{Sampling inputs with a random initial basis}
\label{frames:average-section}

To begin contraction before a deterministic first sweep, choose
$v_{t+1}$ uniformly outside the span $W_t$ of the last
$\min(t,d-1)$ directions, with $W_0=\{0\}$. Use independent
fair switches conditional on the complete schedule. Each
successive set of at most $d$ directions is independent, so
Lemma~\ref{lem:frames} applies when $T\ge d$. Choose a
hyperplane $J_t\supseteq W_t$ by a fixed rule based on the
direction prefix.

Let $C$ be a uniform ordered list, independent of the joint
schedule and process, and take $\Xi=C$. The online energy
calculation uses only the chronological field
\eqref{eq:marginal-filtration}. At terminal time
Lemma~\ref{shear:lem:tuple} first bounds
\[
 \E_{\mathbf v,C}
 \|\mathcal L(Z_T(C)\mid\mathbf v,C)-\pi_q\|_{\TV}.
\]
For each fixed schedule, the general frame comparison then gives
\begin{equation}\label{shear:eq:sampled-frame-transfer}
 \E_C\|\mathcal L(Z_T(C)\mid\mathbf v,C)-\pi_q\|_{\TV}
 =
 \E_C\|\mathcal L(X_T(C)\mid C)-\pi_q\|_{\TV}.
\end{equation}
Here the reference coordinate order is $(1,\ldots,d)$.
The map $L_0^{-1}$ disappears only after averaging $C$. It is
determined once the first $d$ directions have been revealed; we
use it only in the terminal comparison and do not adjoin it
before those defining directions are available.

\begin{proposition}[Averaged tuples with no initial deterministic sweep]
\label{shear:prop:random-tuple}
Let $d\ge3$, $q=7N/8$, and $T=1024d$. For a uniformly sampled
ordered distinct input tuple $C$, independent of $X$,
\begin{equation}\label{shear:eq:random-tuple}
 \E_C\|\mathcal L(X_T(C)\mid C)-\pi_q\|_{\TV}
 \le\frac{q\sqrt N}{2}
 \left(e^{-16d}+2(N+1)e^{-N/256}\right)^{1/2}=o(1).
\end{equation}
The corresponding hidden-card energy is $O(N^{-20})$,
uniformly over its index.
\end{proposition}
\begin{proof}
For the hidden label at index $j$, let $h=N-j+1\ge N/8$.
Conditional on the full direction and switch history but not
on $C$, its available set is a uniform $h$-subset. For its
count in the fixed half $J_t$, choose each point independently
with probability $h/N$ and condition on selecting exactly
$h$ points. That size has probability at least $1/(N+1)$:
it is a mode of the binomial size distribution, as comparison
of consecutive probabilities shows. Before conditioning, the
half-count has mean $h/2$ and is a sum of $N/2$ independent
bounded variables. Lemma~\ref{lem:concentration} gives
\begin{equation}\label{shear:eq:sample-balance}
 \Pr\{\min(|D_t^j\cap J_t|,|D_t^j\setminus J_t|)<h/4\}
 \le2(N+1)e^{-h^2/(4N)}
 \le\beta_N,\qquad
 \beta_N=2(N+1)e^{-N/256}.
\end{equation}
For $h=N$ the count is deterministic and the bound is automatic.

The proper-subspace part of Lemma~\ref{shear:lem:energy}
gives conditional expected loss at least
$h\Delta_t/(8N)\ge\Delta_t/64$ off this event. Hence, from
time zero,
\[
 \E\Delta_{t+1}\le(63/64)\E\Delta_t+\beta_N/64,\qquad
 \E\Delta_T\le e^{-T/64}+\beta_N
             =e^{-16d}+\beta_N=O(N^{-20}).
\]
Apply the terminal tuple lemma with $\Lambda=\mathbf v$ and
then the per-schedule averaged equality
\eqref{shear:eq:sampled-frame-transfer}. This proves
\eqref{shear:eq:random-tuple}.
\end{proof}

\begin{lemma}[Ordinary random-input marginals]
\label{frames:average-law}
Let $d\ge6$, $b=64$, $q=N(1-1/b)$, and $T=2048d$. The
average, over uniform ordered lists of $q$ distinct input
labels, of their time-$T$ physical image-law distance from
$\pi_q$ is $o(1)$.
\end{lemma}
\begin{proof}
Use the random initial basis process above. At a hidden index,
the number $h$ of available positions is at least $N/b$.
Conditional on the complete direction and switch history, and
before sampling the list, this set is a uniform $h$-subset.
Thus \eqref{shear:eq:uniform-subset-moment} bounds its imbalance
across the chosen $J_t$ by $2\vartheta^h$.

On the balanced event the proper-subspace contraction gives
loss at least $\Delta_t/(8b)$. Using the weaker additive term
$2\vartheta^{N/b}$ in the unconditional recurrence yields
\[
 \E\Delta_T\le(1-1/(8b))^T+2T\vartheta^{N/b}.
\]
The terminal tuple comparison and then
\eqref{shear:eq:sampled-frame-transfer} give
\begin{align}
 \E_{\mathbf v,C}
 \|\mathcal L(Z_T(C)\mid\mathbf v,C)-\pi_q\|_{\TV}
 &=\E_C\|\mathcal L(X_T(C)\mid C)-\pi_q\|_{\TV}\notag\\
 &\le\frac12N^{3/2}
 \left((1-1/(8b))^T+2T\vartheta^{N/b}\right)^{1/2}.
 \label{frames:average-tv}
\end{align}
Here $(1-1/(8b))^T\le e^{-4d}$, so the right side tends to
zero. Since $T$ is a multiple of $d$, $X_T=Y_T$ in the
rotating-frame coupling, which proves the physical statement.
\end{proof}
